\documentclass[a4paper,oneside,12pt]{amsart}

\usepackage[spanish]{babel}
\usepackage[utf8]{inputenc}
%\usepackage[latin1]{inputenc}
\usepackage{latexsym}
\usepackage{multicol}
\usepackage{enumerate}
\usepackage{enumitem}
\usepackage[heightrounded]{geometry}
\geometry{top=2cm,bottom=2cm,left=2cm,right=2cm}
\usepackage{graphicx}
\usepackage{tabularx}
\usepackage[spanish]{cleveref}
\usepackage{tikz}
\usetikzlibrary{patterns} % Load the patterns library
\usepackage{caption}
\usepackage{subcaption}
\usepackage{booktabs}

\DeclareMathOperator{\cond}{cond}
\DeclareMathOperator{\Id}{Id}
\DeclareMathOperator{\re}{Re}
%\DeclareMathOperator{\im}{Im}
\newcommand{\I}{\mathbb{I}}
\newcommand{\R}{\mathbb{R}}
\newcommand{\Q}{\mathbb{Q}}
\newcommand{\C}{\mathbb{C}}
\newcommand{\Z}{\mathbb{Z}}
\newcommand{\N}{\mathbb{N}}
\DeclareMathOperator{\im}{Im}

\newcommand{\nombremateria}
{Investigaci\'on Operativa (M) \\ Introducción a la Investigación Operativa y Optimización (LCD)}
\newcommand{\nombrecuatrimestre}
    {Segundo cuatrimestre de 2025}
\newcommand{\numeroparcial}
    { Primer Parcial }
\newcommand{\fecha}
    {5 de Septiembre de 2025}
\newcommand{\numerotema}
    {\bf 1}

%\oddsidemargin -1cm \evensidemargin -1cm \textwidth 17.5cm
%topmargin 1.2cm \textheight 25cm
%\setlength{\textheight}{25cm}

\def \ds{\displaystyle}
\theoremstyle{definition}
\newtheorem{quest}{Ejercicio}
\thispagestyle{empty}

\begin{document}

    \begin{minipage}[t]{0.55\textwidth}
    \begin{tabularx}{\textwidth}{Xl}
        Nombre: & LU: \hphantom{5cm}\\[0.5em]
        Nombre: & LU: \\[0.5em]
        Nombre: & LU: \\[0.5em]
    \end{tabularx}
    \end{minipage}
    \hfill
    \begin{minipage}[t]{0.35\textwidth}
        \begin{tabular}[c]{|c|c|c|}
            \hline
            1 & 2 & Calificación \\ \hline \hline
            \hspace{1.2 cm} & \hspace{1.2 cm} & \hspace{1.5 cm} \\
            \hspace{1.2 cm} & \hspace{1.2 cm} & \hspace{1.5 cm} \\ \hline
        \end{tabular}
    \end{minipage}
%\hfill \hfill

%\vspace{-1cm}
%
%\noindent Nombre y apellido:
%
%\vspace{0.2cm}
%
%\noindent Número de libreta:

%\noindent
\hrulefill

\smallskip
\renewcommand{\baselinestretch}{1.1}

\begin{center}
	\large \textbf{\nombremateria} 
	
	\numeroparcial -- \fecha
\end{center}

\vspace{.6cm}
\renewcommand{\theenumi}{\textbf{\arabic{enumi}}}

\begin{center}
    \emph{Entregar el Ejercicio 1 y el Ejercicio 2 en hojas separadas.}
\end{center}
%%% EMPIEZA EL EJERCICIO 1 %%%

\begin{quest} 
    \texttt{[5 pts.]}En un barrio con forma de grilla de $n\times n$ manzanas, se desean instalar cabinas de
    seguridad en las
    esquinas.
    \begin{enumerate}[label=\alph*)]
        \item \texttt{[2 pts.]} Decimos que una esquina está vigilada si hay una cabina en ella o en alguna esquina
        adyacente. Por
        ejemplo, en la \Cref{fig:1a}, si se instala una cabina en la esquina negra, esa esquina y las esquinas
        sombreadas están vigiladas.
        Elaborar un modelo de Programación Lineal Entera para decidir dónde instalar las cabinas, de manera tal que
        todas las esquinas estén vigiladas y se minimice la cantidad de cabinas instaladas.
        \item \texttt{[1 pt.]}Utilizando SCIP, resolver el modelo para la instancia de una grilla de $4\times 4$
        manzanas. En la hoja, escribir el valor de las variables en la solución obtenida, graficar la grilla y marcar
        dónde se ubican las cabinas de seguridad.
        \item \texttt{[2 pts.]}Decimos que una cuadra está vigilada si al menos en una de sus esquinas hay una
        cabina instalada. Por ejemplo, en la \Cref{fig:1b}, si se instala una cabina en la esquina negra,
        las cuadras punteadas están vigiladas.
        Elaborar un modelo de Programación Lineal Entera para decidir dónde instalar las cabinas, de manera tal que
        todas las cuadras estén vigiladas  y se minimice la cantidad de cabinas instaladas.
        \item \texttt{[1 pt.]}Utilizando SCIP, resolver el modelo para la instancia de una grilla de $4\times 4$
        manzanas. En la hoja, escribir el valor de las variables en la solución obtenida, graficar la grilla y marcar
        dónde se ubican las cabinas de seguridad.
    \end{enumerate}

    \begin{figure}[h!]
        \centering
        \begin{subfigure}[t]{0.48\textwidth}
            \centering
            \begin{tikzpicture}
                \tikzstyle{normal}=[fill=white, draw=black, shape=circle]
                \tikzstyle{marked}=[pattern=north east lines, draw=black, shape=circle]
%                \tikzstyle{marked}=[fill=blue, draw=black, shape=circle]
                \tikzstyle{marked 2}=[fill=black, draw=black, shape=circle]
                \draw[step=1cm] (0,0) grid (4,4);
                \foreach \x in {0, 1, 2, 3, 4}
                    \foreach \y in {0,1,2,3,4}
                        \node[style=normal] () at (\x, \y) {};
                \node[style=marked 2] () at (1,3) {};
                \node[style=marked] () at (1,4) {};
                \node[style=marked] () at (0,3) {};
                \node[style=marked] () at (2,3) {};
                \node[style=marked] () at (1,2) {};

            \end{tikzpicture}
            \caption{Si se instala una cabina en la esquina negra, esa esquina y las esquinas sombreadas están vigiladas.}
            \label{fig:1a}
        \end{subfigure} \hfill
        \begin{subfigure}[t]{0.48\textwidth}
            \centering
            \begin{tikzpicture}
                \tikzstyle{normal}=[fill=white, draw=black, shape=circle]
                \tikzstyle{marked 2}=[fill=black, draw=black, shape=circle]
%                \draw[step=1cm] (0,0) grid (4,4);
                \draw[step=1cm] (2,0) grid (4,4);
                \draw[step=1cm] (0,0) grid (2,2);
                \draw[ultra thick, dashed] (1,3) to (1,2);
                \draw[ultra thick, dashed] (1,3) to (1,4);
                \draw[ultra thick, dashed] (1,3) to (0,3);
                \draw[ultra thick, dashed] (1,3) to (2,3);
                \draw (0,2) to (0,3);
                \draw (0,3) to (0,4);
                \draw (0,4) to (1,4);
                \draw (1,4) to (2,4);
                \foreach \x in {0, 1, 2, 3, 4}
                    \foreach \y in {0,1,2,3,4}
                        \node[style=normal] () at (\x, \y) {};
                \node[style=marked 2] () at (1,3) {};
            \end{tikzpicture}
            \caption{Si se instala una cabina en la esquina negra, las cuadras punteadas están vigiladas.}
            \label{fig:1b}
        \end{subfigure}
    \caption{}
    \end{figure}

\end{quest}

%%% TERMINA EL EJERCICIO 1 %%%

%%% EMPIEZA EL EJERCICIO 2 %%%

\begin{quest}
    \texttt{[5 pts.]} La compañía cervecera Cantares fabrica tres tipos de cerveza: $C_1$, $C_2$ y $C_3$.
    Opcionalmente, puede producir una nueva variedad de cerveza $C_4$. Si elabora esta nueva
    variedad, configurar las máquinas cuesta $\$1750$ semanales.

    El centro de distribución requiere que semanalmente se produzcan al menos cierta cantidad de litros de cada tipo
    de cerveza. Si decide producir $C_4$, se requerián al menos $75$ lts semanales.
    Como la empresa tiene un depósito con capacidad limitada, se puede almacenar cierta cantidad de cada
    tipo de cerveza por semana. Si se excede esa cantidad, la empresa debe alquilar lugar en otro depósito, pagando
    $k_i$ pesos por cada litro excedente de $C_i$. La empresa no desea gastar más de $\$2000$
    en alquiler de espacio para almacenamiento semanalmente.

    La elaboración de cada litro de cada variedad de cerveza requiere cierta cantidad de tiempo.
    Como la empresa cuenta con máquinas que operan en paralelo, dispone de hasta $700$ horas semanales para
    producir cerveza. Finalmente, por razones de marketing, el valor absoluto de la diferencia entre las
    cantidades de $C_2$ y de $C_3$ producidas debe ser mayor o igual que $75$ lts.

    Los datos del problema se encuentran en la siguiente tabla:
    \begin{table}[h!]
    \centering
    \begin{tabular}{lccccc}
         & $C_1$ & $C_2$ & $C_3$ & $C_4$ \\\midrule
        Precio de venta (por lt.) & 8 & 19 & 30 & 52 \\
        Demanda del centro de distribución (en lts.) & 250 & 210 & 150 & 75 \\
        Capacidad de almacenamiento (en lts.) & 275 & 210 & 250 & 0 \\
        Costo de almacenamiento extra (por lt.) & 2 & 5 & 8 & 10 \\
        Horas de elaboración (por lt.) & 0.18 & 0.25 & 0.45 & 0.65 \\\bottomrule
    \end{tabular}
\end{table}

    \begin{enumerate}[label=\alph*)]
        \item \texttt{[3 pts.]} Formular un modelo de Programación Lineal Entera Mixta cuyo objetivo sea maximizar la
        ganancia semanal
        de la empresa, considerando los gastos que podrían implicar la producción de $C_4$ y el alquiler de espacio
        en un depósito.
        Como se vende la cerveza en botellas de distintos tamaños, la cantidad producida de cada variedad puede ser
        una magnitud real.
        \item \texttt{[2 pts.]} Resolver el modelo utilizando SCIP. En la hoja, responder cuántos litros de cada
        tipo de cerveza deben producirse para maximizar la ganancia.
    \end{enumerate}

\end{quest}

%%% TERMINA EL EJERCICIO 2 %%%

%%% TERMINA EL PARCIAL %%%



\end{document}

